\section{Discussion and Qualitative Analysis}
\label{sec:discussion}
The experimental results indicate that PhyST-Net provides a significant leap in grid-ready forecasting. In this section, we delve into the qualitative reasons for this performance shift and discuss the broader implications for renewable energy integration across diverse climatic zones.

\subsection{The Role of Temporal Continuity in Tokenization}
\label{sec:tokenization_role}
Our analysis reveals that the differentiable nature of adaptive Gaussian-soft-windowed patching (AGSP) is far more than a mathematical convenience. By allowing the Gaussian windows to shift their centers ($\mu$) and widths ($\sigma$), the model effectively learns an \textbf{adaptive low-pass filter}. During stable atmospheric periods, the windows expand to perform a robust smoothing of the signal, ignoring sensor jitter and high-frequency turbulent fluctuations that do not contribute to bulk power generation. Conversely, when the model detects an incoming wind ramp, the $\sigma$ parameters contract, allowing the tokens to focus on the high-frequency wavefront of the ramp. This behavior is fundamentally impossible with the rigid patching used in traditional models like PatchTST, suffering from artifacts that distort the gradient signal during rapid transitions.

Furthermore, we observe that the learnable centers $\mu_i$ tend to cluster around periods of high information density. This unsupervised adaptive centering—which we define as ``event alignment''—ensures that the latent tokens provided to the subsequent spatial layers are semantically rich and physically consistent, capturing the true multi-scale nature of the atmospheric signal.

\subsection{Spatial Topology Awareness and Disentanglement}
\label{sec:topology_awareness}
One of the most striking observations from our experiments is the model's ability to ``see'' the farm-wide spatial structure through the relational spatio-temporal manifold fusion (R-STMF) framework. The Far-Field branch of R-STMF learns dense connections between turbines exhibiting high synoptic correlation, even if they are not geographical nearest neighbors. For instance, in the SDWPF dataset, turbines located on the windward side of the farm show a high correlation with those on the leeward side during stable monsoon events.

Meanwhile, the Seasonal branch, with its directional-distance kernel, captures the immediate wake deficits from upwind neighbors. This branch is particularly active during periods of moderate wind speeds where wake effects are most pronounced. By disentangling these spatial scales, PhyST-Net avoids the ``spatial information smearing'' that plagues traditional single-graph graph neural networks (GNNs), often failing to distinguish between a localized turbine issue and a macroscopic weather front moving across the province.

\subsection{K-AMC: Leveraging Spline Physics for NWP Integration}
\label{sec:kamc_physics}
The superior performance of KAN-enhanced adaptive meteorological conditioning (K-AMC) over standard MLP-based modulation validates our hypothesis about the ``KAN Advantage''—specifically, the theoretical capability of KAN spline-based edges to adaptively approximate piecewise continuous mappings without the global gradient smearing typical of traditional MLPs. The mapping from wind speed to power (the power curve) is not a simple linear or quadratic function; it is a piecewise continuous mapping with specific physical ``knees'' at the cut-in and rated speeds. Kolmogorov-Arnold networks (KANs), with their spline-based activations on the edges, are naturally suited to approximate such complex physical shapes. 

Visualization of the K-AMC edge functions (as represented by the localized correction weights in Fig. 10) shows that the model learns to ``attenuate'' the influence of numerical weather prediction (NWP) temperature features once the wind speed exceeds the turbine's rated threshold. This is a sophisticated physical behavior: at high wind speeds, the turbine is governed by pitch control limits rather than air density fluctuations. Traditional multi-layer perceptrons (MLPs) struggle to learn such localized, piecewise behavior without affecting the global weight distribution.

\subsection{Generalizability Across Climatic Zones}
\label{sec:generalizability}
A critical requirement for industrial WPF deployment is generalizability. PhyST-Net was validated on three sites with vastly different physical signatures: coastal (SDWPF), highland (HoT), and small-cluster (Kelmarsh). 

In the coastal environment of Shandong, the model's Trend branch effectively captured the sea-land breeze cycle, driven by synoptic temperature gradients. In the Scottish Highlands, the AGSP module's ability to contract its bandwidth was essential for handling the steep wind ramps induced by rugged terrain. The consistent performance across these diverse regimes suggests that the physical inductive biases embedded in PhyST-Net are universally applicable, reducing the need for extensive site-specific hyperparameter tuning.

\subsection{Societal Impact and Environmental Sustainability}
\label{sec:societal_impact}
The adoption of high-fidelity WPF models like PhyST-Net has profound societal implications. As nations decommission fossil-fuel power plants, the burden of ensuring a reliable electricity supply shifts to renewable assets. Inaccurate forecasts can lead to blackouts or the need for carbon-intensive ``peaker'' plants to be kept on standby. 

By achieving a 96.02\% grid compliance rate, our model supports the stable transition to a net-zero grid. This reduces the ``carbon footprint'' of the balancing market itself, as grid operators can rely more on wind and less on spinning reserves. Furthermore, more accurate forecasts reduce the price volatility of the electricity market, directly benefiting residential and industrial consumers. The intersection of physics-informed AI and renewable energy is thus not just a technical frontier, but a key enabler of global climate justice and energy equity.

\subsection{Grid Stability and Physical Safety via Structural Projection}
\label{sec:stability_safety}
The introduction of the PI-DCH has profound implications for grid frequency stability. In high-penetration scenarios, even a small number of ``physically impossible'' forecasts---such as over-predicting power by 10\% during a storm---can lead to sub-optimal dispatch decisions that trigger protective relays. 

By projecting the output manifold onto the physically feasible domain $[0, P_{rated}]$, PhyST-Net ensures that grid operators receive a signal that is always actionable and physically grounded. The dual-Sigmoid clipper acts as a ``Structural Regularizer,'' preventing the model from following spurious gradients leading to mechanical limit violations. This ``Safety-First'' architecture is essential for the transition from research prototypes to industrial-grade power system components that can be trusted by system operators.

\subsection{Economic Impact of High Compliance Forecasting}
\label{sec:economics}
The achievement of a 96.02\% grid compliance rate is not merely a statistical victory; it has significant economic implications. By reducing the frequency of large-scale forecast errors, PhyST-Net minimizes the need for high-cost spinning reserves.

Our analysis suggests that for a mid-sized 100MW wind farm, the adoption of PhyST-Net could result in an estimated annual saving of over \$250,000 in operational costs compared to standard baseline models (representing a preliminary economic projection based on reduced regulatory penalties and lowered operating reserve requirements). This cost reduction comes from three primary sources: (1) Reduced Balancing Costs: Fewer deviations between day-ahead schedules and actual generation; (2) Lower Penalty Payments: Improved adherence to grid code requirements for frequency response; (3) Asset Longevity: Physically bounded forecasts lead to smoother dispatch commands, reducing the mechanical wear and tear on pitch systems and drive trains during volatile weather.
